\begin{pastexam}[label={px:ust-107-ee-ds-07}]{臺聯大電機類 107 學年度 第七題}
（問答；10 分）

(10\%) Answer the following questions about a certain type of data structures.
\begin{enumerate}
\item[(一)] (2\%) Which data structure is often used to capture relationships among entities?
\item[(二)] (4\%) Entities are represented as \blank{1} and relationships are represented as \blank{2}.
\item[(三)] (4\%) Depending on whether or not the relationships are symmetric, the data structure can be divided into two subtypes \blank{3} and \blank{4}.
\end{enumerate}
\end{pastexam}

\begin{pastsoln}
（空格的編號 (1)--(4) 是本書加的，方便對答案。）

(一)~圖（graph）。

(二)~(1) 頂點（vertices，也叫 nodes）；(2) 邊（edges）。 圖 $G = (V, E)$，$V$ 是頂點集合，每個實體是一個頂點；
$E$ 是邊的集合，兩個實體之間有關係，就在對應的兩個頂點之間連一條邊（資結書第~21.1~節）。

(三)~(3) 無向圖（undirected graph）；(4) 有向圖（directed graph，digraph）。 關係是對稱的（$u$ 和 $v$ 有關係就表示
$v$ 和 $u$ 有關係，例如「是朋友」「有道路相連」），邊沒有方向，用無序對 $\{u, v\}$ 表示，是無向圖；關係不一定對稱
（例如「追蹤」「單行道」「是先修課」），邊有方向，用有序對 $(u, v)$ 表示，是有向圖。 無向圖的相鄰矩陣是對稱的，
有向圖則不一定。
\end{pastsoln}
